\originalchapter{超导响应、相干因子与弱连接}
\label{ch:scresponse}
\source{8.512第9--13讲及8.511配对基础。}
\originalsection{横向响应与极限顺序}
正常电子气对静态、空间均匀的纯规范$A$不能产生可观测电流。其顺磁与抗磁响应在$\omega=0,q\to0$抵消；Landau轨道抗磁来自保留空间变化至适当阶的横向响应，不能把常数$A$响应当作磁化。超导体在横向静态极限保留刚度$K_T(0,0)=n_se^2/m$，于是$\vct j_T=-K_T\vct A_T$，恢复London方程。

无散射正常金属在先$q=0$后$\omega\to0$的动态极限也有Drude零频权重。理想电导与Meissner排磁区别在于不同极限的响应，不是只看$\sigma(0)$是否无穷大。规范不变计算还应让序参量相位随纵向扰动调整，否则会违反电荷守恒。
\originalsection{热准粒子与超流密度}
在各向同性洁净$s$波体系，流动使准粒子能量有Doppler位移。准粒子对动量的反向响应给
\begin{equation}
\frac{n_s(T)}n=1+2\int_\Delta^\infty\frac{E\dd E}{\sqrt{E^2-\Delta^2}}\frac{\partial f(E)}{\partial E}.
\end{equation}
零温$\partial f/\partial E$在正能量区为零，$n_s=n$。低温令$E=\Delta+x$，平方根分母$\sqrt{2\Delta x}$，$-f'(E)\simeq\e^{-E/\kb T}/\kb T$，积分给$n-n_s\propto n\sqrt{2\pi\Delta/\kb T}\e^{-\Delta/\kb T}$。穿透深度$\lambda^{-2}(T)=\mu_0n_s(T)e^2/m$。这一低温指数律又依赖无节点能隙。

非磁性弹性杂质会破坏单个$k$态，却不必破坏由时间反演态构成的$s$波配对。若杂质不改变平均态密度且配对近各向同性，能隙方程中频率与能隙的重整化有相同因子，$T_c$在该近似内保持。这是Anderson定理的内容，不是所有无序都无影响。超流刚度与穿透深度仍会变化；脏极限$\ell\ll\xi_0$有$\xi_{\rm dirty}\sim\sqrt{\xi_0\ell}$，$\lambda_{\rm dirty}\sim\lambda_L\sqrt{\xi_0/\ell}$。磁性散射翻转自旋或非$s$波配对中的角向混合通常破坏这项保护。
\originalsection{相干因子的来源}
把电子算符展开为Bogoliubov准粒子，一般双线性观测量包含准粒子散射系数$u_ku_{k'}\mp v_kv_{k'}$和两准粒子创建系数$u_kv_{k'}\pm v_ku_{k'}$。符号由观测量在时间反演下的性质和自旋结构决定，不是任意挑选。

例如密度与自旋观测量的电子、空穴贡献可相消或相加。在费米面附近$u\simeq v$，某些通道被压低，另一些通道增强。因此超导隙上方态密度的平方根峰并不保证所有散射实验都有同样强的峰。核磁弛豫中的相干峰需要特定矩阵元及散射条件；相位改变的配对、寿命展宽与强涨落可削弱它。

做计算时应先写待测算符，再代入$u,v$，再按准粒子初末态取矩阵元。若只写一个未说明适用通道的“相干因子”，难以判断符号是否正确。
\originalsection{隧穿电流与Josephson微观结构}
正常金属--超导体的弱隧穿，在矩阵元近常数时$dI/\dd V\propto D_s(eV)$，低温可测能隙与谱峰。超导--超导体的耗散电流有两边态密度卷积，零温阈值$|eV|\ge\Delta_1+\Delta_2$。配对电流来自隧穿哈密顿量二阶过程：先转移一个电子的中间态有准粒子能量代价，第二个转移恢复配对，两种方向振幅干涉产生$\cos\varphi$能量。

相同$s$波隙、弱隧穿且无磁性时，微观能量积分给Ambegaokar--Baratoff关系
\begin{equation}
I_cR_N=\frac{\pi\Delta(T)}{2e}\tanh\frac{\Delta(T)}{2\kb T}.
\end{equation}
$R_N$是同一结的正常电阻，公式不能用于任意高透明度接触。在超导环的两个结中，相位差相差$2\pi\Phi/\Phi_0$；相同结总临界流$2I_c|\cos(\pi\Phi/\Phi_0)|$，若忽略环自感。这把规范相位与可测磁通直接联系起来。
